\chapter{Dyspraxia}
\index{Dyspraxia}

\section*{Definition and Overview}
Dyspraxia, also known as developmental coordination disorder (DCD), is a condition that affects the planning and coordination of physical movements. Children and adults with dyspraxia have difficulty with motor skills that are out of proportion to their age and intelligence, such as tying shoelaces, riding a bike, or handwriting, and these difficulties interfere with daily life and learning. Dyspraxia is a lifelong condition with a neurological basis, but with early recognition, appropriate therapy, and support, most people learn to manage their difficulties and develop their strengths.

\section*{Epidemiology}
Dyspraxia affects an estimated 5 to 6 per cent of school-aged children and is more common in boys. It frequently coexists with other conditions, including attention deficit hyperactivity disorder, autism, dyslexia, and speech and language difficulties. Many children with dyspraxia are not identified until they reach school age, when the demands of handwriting, sport, and self-care expose their motor difficulties.

\section*{Aetiology and Pathophysiology}
The cause of dyspraxia is not fully understood, but it reflects differences in how the brain plans and executes movements.
\begin{itemize}
    \item \textbf{Brain function:} difficulties in the brain's ability to process sensory information and plan motor actions, without any damage to the muscles or nerves themselves
    \item \textbf{Genetic factors:} dyspraxia tends to run in families
    \item \textbf{Prematurity and birth factors:} low birth weight and premature birth are associated with higher rates
    \item \textbf{Coexisting conditions:} ADHD, autism, and language disorders commonly co-occur
    \item \textbf{Sensory processing:} many affected children have differences in processing touch, movement, and spatial information
    \item \textbf{The motor learning cycle:} each movement relies on planning, executing, and learning from feedback, all of which can be impaired
\end{itemize}

\section*{Clinical Features}
The features of dyspraxia change across development.
\begin{itemize}
    \item Delayed motor milestones, such as sitting, crawling, and walking
    \item Clumsiness, frequent bumping into things, and poor balance
    \item Difficulty with fine motor tasks: handwriting, using cutlery, doing up buttons, and tying shoelaces
    \item Difficulty with gross motor skills: running, jumping, catching, and riding a bike
    \item Awkward or unusual gait and posture
    \item Difficulty learning new movements, which may improve with practice but slowly
    \item Speech and language difficulties, including articulation problems
    \item Difficulty with organisation, planning, and time management
    \item Fatigue and avoidance of physical activity
    \item In adults: difficulty with driving, DIY, and new physical skills
\end{itemize}

\section*{Differential Diagnosis}
Other conditions can cause motor difficulties and must be excluded.
\begin{itemize}
    \item Muscular conditions such as muscular dystrophy, which cause progressive weakness
    \item Neurological conditions including cerebral palsy
    \item Intellectual disability or global developmental delay
    \item Visual or hearing impairment affecting coordination
    \item Attention deficit hyperactivity disorder, which affects attention rather than coordination itself
    \item Anxiety, which can impair performance of motor tasks
\end{itemize}

\section*{When to Seek Medical Care}
Parents and teachers should seek assessment when a child's motor skills are significantly behind expectations or are causing distress, frustration, or avoidance of activity. A GP can coordinate referral to a paediatrician, occupational therapist, or physiotherapist.

\textbf{Contact your local emergency services for:}
\begin{itemize}
    \item Sudden loss of coordination or new weakness in an adult, which may indicate a stroke
    \item Sudden confusion, difficulty speaking, or facial drooping
\end{itemize}

\section*{Diagnosis and Investigations}
Diagnosis of developmental coordination disorder is made by a multidisciplinary team.
\begin{itemize}
    \item \textbf{History:} developmental milestones, motor skills, and impact on daily life
    \item \textbf{Standardised motor assessment:} tests such as the Movement ABC measure motor proficiency against age norms
    \item \textbf{Cognitive assessment:} to confirm that motor difficulties are not explained by intellectual disability
    \item \textbf{Occupational therapy assessment:} evaluation of daily living skills and fine motor function
    \item \textbf{Medical examination:} to exclude neurological, muscular, and sensory causes
    \item \textbf{Screening for coexisting conditions:} assessment for ADHD, autism, and language disorders
\end{itemize}

\section*{Management}
Management is supportive and therapeutic, focused on function and confidence.
\begin{itemize}
    \item \textbf{Occupational therapy:} task-based training to break down movements and practise them systematically
    \item \textbf{Physiotherapy:} exercises for balance, strength, and gross motor skills
    \item \textbf{School support:} accommodations such as extra time for motor tasks, keyboard use, and reduced copying
    \item \textbf{Task-specific training:} practising the specific skills the child needs, such as shoelaces or bike riding
    \item \textbf{Adaptive equipment:} pencil grips, weighted cutlery, and other aids
    \item \textbf{Speech therapy:} for associated speech and language difficulties
    \item \textbf{Emotional support:} building self-esteem and addressing frustration and anxiety
    \item \textbf{Adult support:} workplace adjustments and strategies for driving and daily tasks
\end{itemize}

\section*{Complications}
\begin{itemize}
    \item Low self-esteem, frustration, and avoidance of physical activity
    \item Bullying and social difficulties at school
    \item Academic underachievement related to handwriting and organisational difficulties
    \item Anxiety and depression
    \item Reduced participation in sport and social activities
    \item Fatigue from the extra effort required for everyday tasks
\end{itemize}

\section*{Prognosis and Outcomes}
Dyspraxia is a lifelong condition, but motor skills usually improve with age, practice, and targeted therapy, and many adults with dyspraxia manage well with strategies and accommodations. The condition does not affect intelligence, and many people with dyspraxia are creative, determined, and successful. Early identification, supportive teaching, and attention to emotional wellbeing are the keys to good outcomes.

\section*{Prevention and Self-Care}
\begin{itemize}
    \item Seek assessment early to access therapy and support
    \item Practise specific skills in small, achievable steps
    \item Allow extra time for motor tasks and reduce pressure
    \item Use keyboarding and assistive technology for writing
    \item Encourage physical activity in supportive, non-competitive settings
    \item Celebrate effort and strengths alongside the difficulties
    \item Ensure adequate rest, as motor tasks are more tiring
\end{itemize}

\section*{Sources and Further Reading}
The following reputable sources provide further information for healthcare students and patients seeking reliable, current advice about dyspraxia.
\begin{itemize}
    \item Dyspraxia Foundation. \textit{Dyspraxia information and support.} https://dyspraxiafoundation.org.uk/
    \item Occupational Therapy Australia. \textit{Developmental coordination disorder resources.} https://www.otaus.com.au/
    \item SPELD NSW and other state associations. \textit{Developmental coordination disorder information.} https://www.speldnsw.org.au/
    \item NHS. \textit{Developmental coordination disorder (dyspraxia).} https://www.nhs.uk/conditions/developmental-coordination-disorder-dyspraxia/
    \item Mayo Clinic. \textit{Developmental coordination disorder.} https://www.mayoclinic.org/diseases-conditions/developmental-coordination-disorder/
    \item Understood.org. \textit{Dyspraxia: what it is and how to help.} https://www.understood.org/
\end{itemize}
